\chapter{Immersed Boundary Methods}
\label{chap:immersed-boundary-methods}

\section{Introduction}

\emph{Fluid--structure interaction} (FSI) problems couple the motion of a fluid to the
motion or deformation of an immersed solid. The two subsystems are not
independent: the fluid exerts hydrodynamic traction on the structure, while the
structure changes the admissible fluid velocity through no-slip,
no-penetration, or more general interfacial conditions. Numerical FSI methods
are therefore usually distinguished less by the governing equations themselves
than by how the interface is represented and how strongly the fluid and
structural subproblems are coupled during each time step.

The physical coupling may be modeled as a one-way or two-way coupling. In a
one-way calculation, the structure is prescribed or is otherwise advanced using
fluid loads without allowing its motion to feed back onto the fluid. This is
useful for imposed kinematics or post-processing force estimates, but it is not
a fully coupled FSI problem. By contrast, in two-way FSI, the fluid and
structure mutually determine one another; fluid tractions drive the structural
motion, and the resulting structural configuration changes the fluid domain,
boundary condition, or body-force field. The remaining classifications in this
section mainly concern how this two-way coupling is represented numerically.

One broad distinction is between \emph{boundary-fitted} and
\emph{non-boundary-fitted} methods. In a boundary-fitted formulation, the fluid
mesh conforms to the fluid--solid interface. Arbitrary Lagrangian--Eulerian
(ALE) methods represent the standard example: In this method, the structure is
represented in Lagrangian coordinates, while the fluid mesh moves or deforms to
track the interface. This gives a sharp representation of the boundary
conditions and hydrodynamic tractions, but mesh quality becomes a central issue
when the structure undergoes large displacements, rotations, contact, or
topology changes. In contrast, a large body of non-boundary-fitted methods keep
the fluid grid independent of the solid geometry. Immersed boundary, immersed
interface, fictitious-domain, cut-cell, and related embedded-boundary methods
all fall into this class. These methods avoid remeshing and are particularly
attractive for large motion or many particles, at the cost of introducing
additional machinery to transfer velocity, force, or constraint information
between the Eulerian fluid grid and the Lagrangian solid description.

FSI methods are also commonly categorized by their coupling architecture
\cite{hou_numerical_2012}. In a \emph{monolithic method}, the fluid unknowns,
structural unknowns, and interface constraints are assembled into a single
nonlinear algebraic system at each time step. This formulation treats the
coupled problem as one system, which can be advantageous for stability and for
problems with strong added-mass effects, but it requires solvers and
preconditioners for a large multiphysics system. In a \emph{partitioned
method}, existing fluid and structural solvers are retained as separate modules
and communicate through interface data. The fluid solver may provide tractions
or body forces to the structure, while the structural solver returns
displacements, velocities, or constraint targets to the fluid. This modularity
is attractive in practice, especially when mature solvers already exist for
each subsystem, although the coupling algorithm must be designed carefully so
that splitting errors do not dominate the dynamics, or otherwise accept these
splitting errors as a trade-off in exchange for gained modularity.

Within partitioned methods, a further distinction can be made between weak and
strong coupling. A weakly coupled or explicit scheme advances each subsystem
once per time step using interface data extrapolated from the previous time
level or from an intermediate predictor. Such schemes are simple and cheap, but
they can suffer from severe stability restrictions when the structural inertia
is small compared with the displaced fluid inertia. This is the classical
added-mass difficulty. A strongly coupled or implicit partitioned scheme
instead iterates between the fluid and structure within a time step until the
interface velocity, displacement, and traction conditions are satisfied to a
chosen tolerance. Strong coupling reduces splitting error and is usually more
robust for light or flexible structures, but it increases the cost of each time
step and requires convergence criteria, relaxation, or quasi-Newton
acceleration.

Another useful taxonomy concerns how the interface conditions are imposed. Some
methods impose them strongly by making the fluid and structure share the same
interface degrees of freedom, as in many boundary-fitted finite-element
formulations. Other methods impose the interface equations in a variational or
algebraic sense, for example through Nitsche terms, penalty forces, feedback
laws, or Lagrange multipliers. In immersed boundary methods, the coupling is
commonly expressed through spreading and interpolation operators: a Lagrangian
force density is spread to the Eulerian fluid equations, and the Eulerian
velocity is interpolated back to the Lagrangian markers. Direct-forcing
variants choose the force so that the marker velocity approximately satisfies
the desired no-slip constraint after the fluid
update\cite{uhlmann_immersed_2005}. For flexible filaments or beams, the same
framework can be combined with a structural evolution equation and additional
internal constraints, such as inextensibility\cite{huang_simulation_2007}. From
this viewpoint, the immersed-boundary force plays the role of an interfacial
constraint force: it is not simply an external load, but the numerical
mechanism by which the fluid and structure are made kinematically compatible.

These classifications are not mutually exclusive. A single FSI algorithm may be
non-boundary-fitted, partitioned, weakly coupled in time, and based on a
penalty or direct-forcing approximation of the no-slip constraint. Another may
be boundary-fitted, monolithic, fully implicit, and impose interface continuity
strongly. The choice is therefore problem dependent. Boundary-fitted and
monolithic approaches offer a sharp and tightly coupled description of the
interface, while immersed and partitioned approaches trade some interface
sharpness or implicitness for algorithmic simplicity, robustness under large
motion, and compatibility with fixed-grid flow solvers. The direct-forcing
immersed-boundary methods considered in this work belong to the latter family:
the fluid is solved on an Eulerian grid, the structure or immersed surface is
represented by Lagrangian markers, and the coupling force is constructed to
reduce the slip between the two descriptions.

\section{Immersed Boundary Methods}

Peskin's original immersed boundary method was developed for problems in which
thin elastic structures move inside an incompressible viscous fluid, with heart
valves as the motivating example\cite{peskin_flow_1972,peskin_immersed_2002}.

The central idea is to avoid constructing a body-fitted mesh around the moving
structure. Instead, the fluid variables are stored on a fixed Eulerian grid,
while the immersed structure is represented by Lagrangian material points. The
two descriptions communicate through a regularized delta function. Forces
computed on the Lagrangian structure are spread to the surrounding fluid grid,
and the Eulerian fluid velocity is interpolated back to the Lagrangian
structure to update its motion.

A discrete approximation to dirac-delta function is chosen for the operator
both in spreading and interpolation, acting as the discrete transfer kernel
between the Eulerian and Lagrangian descriptions. There are several common
Peskin-type kernels, differing mainly in support width and smoothness, but they
are designed to satisfy the same basic properties: A constant Eulerian field
should interpolate to the same constant value at a marker, so the kernel
weights should form a partition of unity. The interpolation should not
introduce a systematic shift of the marker position, which motivates a discrete
first-moment condition. The kernel should also be compactly supported over only
a few grid cells so that each marker couples only to a local stencil, while
remaining sufficiently smooth that the force applied to the grid does not jump
violently as a marker crosses cell boundaries.

Equally important, the same kernel is used for interpolation and spreading.
With the appropriate quadrature weights and cell-volume normalization, this
choice makes the two transfer operations adjoint in a discrete virtual-work
sense: the work of a Lagrangian force against the interpolated marker velocity
equals the work of the spread Eulerian force against the Eulerian velocity
field. This property is one reason Peskin-type kernels appear across many
immersed-boundary variants, even when the force itself is computed in very
different ways. Elastic immersed-boundary methods, direct forcing, multi-direct
forcing, and Lagrange-multiplier formulations can all share essentially the
same transfer machinery; they differ primarily in how the Lagrangian force is
chosen.

In the classic IBM formulation, the immersed boundary force is usually
determined from a constitutive law for the structure. For example, if the
boundary is an elastic fiber or membrane, its stretching, bending, or tethering
energy defines a Lagrangian force density. This force is then applied to the
Navier--Stokes equations as a singular body force, regularized over a few grid
cells in the discrete method. The no-slip condition is therefore not imposed by
explicitly solving for a constraint force but rather, implicitly. This is
because the structure is moved with the local fluid velocity (thereby
satisfying no-slip) where after, the elastic force generated by the structural
model feeds back onto the fluid. This gives a compact and robust framework for
moving deformable boundaries without remeshing.

This viewpoint is important because it separates two ingredients that later
methods modify independently: the Eulerian--Lagrangian transfer operators and
the rule used to choose the interaction force. In Peskin's method, the force is
primarily a physical structural force. In direct-forcing immersed boundary
methods, the force is instead chosen so that the interpolated velocity at the
immersed markers satisfies a prescribed boundary velocity after the fluid
update\cite{uhlmann_immersed_2005,fadlun_combined_2000}.

Thus the immersed force becomes closer to a discrete constraint force enforcing
no-slip. Multi-direct forcing extends this idea by repeatedly applying the
interpolation, force correction, and spreading operations within a time step,
so that the residual slip left by a single direct-forcing pass is reduced
\cite{wang_combined_2008}.

In this sense, direct forcing and multi-direct forcing can be viewed as
descendants of Peskin's Eulerian--Lagrangian coupling framework, with a
different choice of how the coupling force is computed.

\section{Distributed-Lagrange Multiplier}

\begin{figure}[H]
    \centering
    \import{figure/vector}{acta-numerica.pdf_tex}
    \caption{}
    \label{fig:immersed-boundary}
\end{figure}

The distributed-Lagrange multiplier/fictitious-domain viewpoint rewrites
fluid--solid coupling as a constraint problem. The fluid equations are solved
on a fixed Cartesian domain that may include the region occupied by the solid,
and a Lagrange multiplier is introduced to enforce the prescribed motion of the
immersed material. In the classical fictitious-domain setting, this multiplier
is used to make the fictitious fluid inside a rigid body move with the rigid
body, rather than only enforcing no-slip at the fluid--solid interface. The
important idea for the present work is more general: the coupling force can be
viewed as a constraint force, analogous to the role pressure plays in enforcing
incompressibility.

The fibres considered here are represented by Lagrangian markers on a space
curve, rather than by a volumetric solid domain. We therefore use the same
constraint-force interpretation, but place the multiplier directly in marker
space. The regularized delta function supplies the interpolation and spreading
operators between the Eulerian grid and the Lagrangian markers, while the
multiplier is stored as a force density on those markers. This is closer in
spirit to immersed-boundary projection methods, such as
\textcite{taira_immersed_2007}, where the boundary force is also interpreted as
a Lagrange multiplier enforcing no-slip at the markers. The main difference is
that Taira and Colonius solve for the boundary force as part of a coupled
projection, so that the projected velocity satisfies both incompressibility and
no-slip. In the present split formulation, the marker force enforces no-slip at
the intermediate velocity before the subsequent pressure projection.

Let \(\vec X_a(t)\), \(a=1,\ldots,N_L\), denote the marker positions and let
\(\vec U_a(t)\) be their prescribed or structural velocity. At the continuous
level, the regularized immersed-boundary coupling can be written schematically
as
\begin{align}
    \rho_f\left(\pdv{\vec u}{t}+\vec u\cdot\nabla\vec u\right)
                      & = -\nabla p + \mu_f\nabla^2\vec u
    -\int_{\Gamma_h}\bm\Lambda(\vec X,t)\,
    \delta_h(\vec x-\vec X)\odif{\Gamma},
    \label{eq:dlmfd-continuous-momentum}                  \\
    \nabla\cdot\vec u & = 0,
    \label{eq:dlmfd-continuous-incompressibility}         \\
    \int_{\Omega}\vec u(\vec x,t)\,
    \delta_h(\vec x-\vec X_a)\odif{\vec x}
                      & = \vec U_a(t),
    \qquad a=1,\ldots,N_L .
    \label{eq:dlmfd-continuous-constraint}
\end{align}
Here \(\delta_h\) is the same regularized kernel used for immersed-boundary
interpolation and spreading. The sign convention in
\eqref{eq:dlmfd-continuous-momentum} is chosen so that
\(\bm\Lambda\) is the force density exerted by the fluid on the marker, while
\(-\bm\Lambda\) is spread back as the force density exerted by the marker on
the fluid. The last equation is the no-slip constraint written at the marker
positions. Thus the interaction force is determined by a constraint equation,
not by a penalty stiffness or a prescribed feedback law.

The discrete method is most easily described in terms of the interpolation and
spreading matrices. Let \(\vec u_i\) be the Eulerian velocity at cell \(i\),
\(\Delta V_i\) the cell volume, \(q_a\) the quadrature weight associated with
marker \(a\), and \(W_{ai}\) the dimensionless kernel weight between marker
\(a\) and cell \(i\). The interpolation operator \(J\) gathers an Eulerian
velocity to the marker locations,
\begin{equation}
    (J\vec u)_a = \sum_i W_{ai}\vec u_i .
    \label{eq:dlmfd-gather}
\end{equation}
The corresponding spreading operator \(S\) maps a marker force density to an
Eulerian force density,
\begin{equation}
    (S\bm\Lambda)_i
    = \frac{1}{\Delta V_i}
    \sum_{a=1}^{N_L} W_{ai} q_a \bm\Lambda_a .
    \label{eq:dlmfd-spread}
\end{equation}
With the Eulerian and marker inner products
\begin{equation}
    (\vec v,\vec u)_\Omega
    = \sum_i \Delta V_i\,\vec v_i\cdot\vec u_i,
    \qquad
    (\bm \eta,\bm\Lambda)_\Gamma
    = \sum_a q_a\,\bm \eta_a\cdot\bm\Lambda_a,
    \label{eq:dlmfd-inner-products}
\end{equation}
The scaling in the spreading operator is motivated by requiring the
Eulerian--Lagrangian transfer to preserve virtual work. The power of the
spread force against any Eulerian velocity field should be the same as the
power of the marker force against the interpolated marker velocity:
\begin{align}
    \sum_i \Delta V_i (S\bm\Lambda)_i\cdot\vec u_i
     & = \sum_a q_a\bm\Lambda_a\cdot(J\vec u)_a \notag \\
     & = \sum_a q_a\bm\Lambda_a\cdot
    \left(\sum_i W_{ai}\vec u_i\right) \notag          \\
     & = \sum_i \vec u_i\cdot
    \left(\sum_a W_{ai}q_a\bm\Lambda_a\right).
    \label{eq:dlmfd-discrete-virtual-work}
\end{align}
Since this relation must hold for arbitrary \(\vec u_i\), comparison with the
Eulerian inner product gives exactly the spread definition in
\eqref{eq:dlmfd-spread}. Thus the factor \(q_a\) appears because the marker
force density is integrated with a Lagrangian quadrature weight, while the
factor \(1/\Delta V_i\) appears because the result is an Eulerian force
density rather than a cell-integrated force.
With this scaling, the two transfer operators are adjoint in the discrete sense,
\begin{equation}
    (S\bm\Lambda,\vec u)_\Omega
    = (\bm\Lambda,J\vec u)_\Gamma .
    \label{eq:dlmfd-adjoint}
\end{equation}
This adjoint relation also fixes the units: \(\bm\Lambda_a\) is a Lagrangian
force density, while \(S\bm\Lambda\) is an Eulerian force density.

Suppose that all non-constraint terms have produced a provisional velocity
\(\vec u^\star\) during a time step. The marker-space DLMFD correction seeks a
velocity correction \(\delta\vec u\) and a force density \(\bm\Lambda\) such
that
\begin{equation}
    \vec u^{n+1}
    = \vec u^\star + \delta\vec u
    \label{eq:dlmfd-velocity-correction}
\end{equation}
satisfies the marker constraint
\begin{equation}
    J\vec u^{n+1} = \vec U^{n+1}.
    \label{eq:dlmfd-discrete-constraint}
\end{equation}
The correction step can be written as the formal coupled system
\begin{equation}
    \begin{bmatrix}
        \mat M_f & \Delta t\,S \\
        J        & 0
    \end{bmatrix}
    \begin{bmatrix}
        \delta\vec u \\
        \bm\Lambda
    \end{bmatrix}
    =
    \begin{bmatrix}
        0 \\
        \vec U^{n+1}-J\vec u^\star
    \end{bmatrix}.
    \label{eq:dlmfd-coupled-correction-system}
\end{equation}
The first row says that the marker force, once spread to the grid, produces the
Eulerian velocity correction. The second row says that the corrected velocity
must satisfy no-slip at the markers. This block system is a useful algebraic
description of the correction step; it is not assembled explicitly in the
implementation.

Here \(\mat M_f\) denotes the diagonal Eulerian fluid mass operator. Since the
spreading operator in \eqref{eq:dlmfd-spread} produces an Eulerian force
density, applying \(\mat M_f^{-1}\) corresponds locally to division by the
fluid density,
\begin{equation}
    \left(\mat M_f^{-1}\vec f\right)_i
    \simeq \frac{\vec f_i}{\rho_i}.
    \label{eq:dlmfd-fluid-mass-inverse}
\end{equation}
For a constant-density calculation, this is simply multiplication by
\(1/\rho_f\). If the finite-volume equations are instead written in integral
form, the cell volume is part of the mass matrix; in the notation used here the
cell-volume factor has already been accounted for in the spreading operator.

Eliminating \(\delta\vec u\) from \eqref{eq:dlmfd-coupled-correction-system}
gives
\begin{equation}
    \delta\vec u
    =
    -\Delta t\,\mat M_f^{-1}S\bm\Lambda,
    \label{eq:dlmfd-eliminated-velocity-correction}
\end{equation}
and substitution into the marker constraint gives the marker-space system
\begin{equation}
    \mat A\bm\Lambda = \vec b,
    \qquad
    \mat A = \Delta t\,J\mat M_f^{-1}S,
    \qquad
    \vec b = J\vec u^\star-\vec U^{n+1}.
    \label{eq:dlmfd-marker-system}
\end{equation}
This reduced marker operator is the Schur complement of the Eulerian velocity
block \(\mat M_f\) in \eqref{eq:dlmfd-coupled-correction-system}, up to the
chosen sign convention. In practical terms, it is the marker mobility: apply a
marker force, spread it to the fluid grid, convert it to a velocity increment,
and measure the result back at the markers.
The right-hand side is the slip velocity left by the provisional fluid update.
Solving \eqref{eq:dlmfd-marker-system} therefore finds the force density that
removes this slip, up to the tolerance of the linear solve and the accuracy of
the regularized transfer operators.

The matrix \(\mat A\) is not assembled explicitly. A matrix-vector product
\(\mat A\bm w\) is evaluated by spreading the trial marker field, applying the
Eulerian mass inverse, and gathering the resulting velocity increment back to
the markers:
\begin{equation}
    \bm y = \mat A\bm w
    = \Delta t\,J\mat M_f^{-1}S\bm w .
    \label{eq:dlmfd-matvec}
\end{equation}
This is the sequence implemented in the code as a spread to a temporary
Eulerian force, division by the local density, and a gather to the temporary
marker vector. Since \(\mat A\) is symmetric positive semi-definite with
respect to the marker inner product in \eqref{eq:dlmfd-inner-products}, the
force density can be computed by conjugate gradients on the marker degrees of
freedom:
\begin{align}
    \alpha_k         & =
    \frac{(\bm r_k,\bm r_k)_\Gamma}
    {(\bm w_k,\mat A\bm w_k)_\Gamma},
                     &
    \bm\Lambda_{k+1} & =
    \bm\Lambda_k+\alpha_k\bm w_k,
    \label{eq:dlmfd-cg-alpha} \\
    \bm r_{k+1}      & =
    \bm r_k-\alpha_k\mat A\bm w_k,
                     &
    \beta_k          & =
    \frac{(\bm r_{k+1},\bm r_{k+1})_\Gamma}
    {(\bm r_k,\bm r_k)_\Gamma},
    \label{eq:dlmfd-cg-beta}  \\
    \bm w_{k+1}      & =
    \bm r_{k+1}+\beta_k\bm w_k .
    \label{eq:dlmfd-cg-direction}
\end{align}
After convergence, the Eulerian IBM forcing applied to the fluid is
\(-S\bm\Lambda\), while the equal and opposite hydrodynamic force on the
immersed object is obtained by summing the marker force densities with their
quadrature weights,
\begin{equation}
    \vec F_h = \sum_{a=1}^{N_L} q_a\bm\Lambda_a .
    \label{eq:dlmfd-hydrodynamic-force}
\end{equation}

The distinction from direct forcing is therefore algebraic. A single
direct-forcing step estimates the marker force locally from the slip velocity,
and multi-direct forcing improves that estimate by repeated corrections. The
marker-space DLMFD method instead solves the coupled marker system
\eqref{eq:dlmfd-marker-system}, so the force at one marker accounts for the
velocity induced at neighbouring markers through the regularized kernel and the
fluid mass operator. The method remains compatible with the same
Eulerian--Lagrangian transfer machinery as the immersed-boundary method, but
the no-slip condition is imposed as a discrete Lagrange-multiplier constraint.

\section{Multi-Direct Forcing Immersed Boundary Method}

The multi-direct forcing (MDF) method uses the same Eulerian--Lagrangian
transfer operators as the marker-space DLMFD method, but it avoids solving the
coupled marker system \eqref{eq:dlmfd-marker-system}. Instead, it builds an
explicit approximation to the inverse of the marker mobility and applies this
correction several times in one time step. In this sense, MDF can be viewed as
a Richardson iteration for the DLMFD Schur complement, preconditioned by a
local pseudoinverse at each marker.

Let \(\vec u^{(k)}\) denote the working Eulerian velocity after \(k\) direct
forcing corrections within a time step, with \(\vec u^{(0)}=\vec u^\star\). The
marker slip is
\begin{equation}
    \vec s^{(k)}
    = \vec U^{n+1} - J\vec u^{(k)} .
    \label{eq:mdf-slip}
\end{equation}
An exact DLMFD correction would choose a marker force-density increment
\(\delta\bm\Lambda^{(k)}\) such that
\begin{equation}
    \vec u^{(k+1)}
    = \vec u^{(k)}
    - \Delta t\,\mat M_f^{-1}S\delta\bm\Lambda^{(k)}
    \label{eq:mdf-velocity-update}
\end{equation}
satisfies \(J\vec u^{(k+1)}=\vec U^{n+1}\). Substituting
\eqref{eq:mdf-velocity-update} into the constraint gives
\begin{equation}
    \Delta t\,J\mat M_f^{-1}S\delta\bm\Lambda^{(k)}
    = -\vec s^{(k)} .
    \label{eq:mdf-exact-increment}
\end{equation}
Thus the exact increment is the solution of the same marker-space operator
\(\mat A=\Delta t\,J\mat M_f^{-1}S\) that appears in
\eqref{eq:dlmfd-marker-system}. MDF replaces this coupled inverse with a local
diagonal approximation.

To derive the local approximation, consider one marker \(a\) and one velocity
component. Let \(I_a\) be the Eulerian cells in the support of its regularized
delta kernel, and collect the corresponding interpolation weights into the row
vector
\begin{equation}
    \mat J_a =
    \begin{bmatrix}
        W_{ai_1} & W_{ai_2} & \cdots & W_{ai_m}
    \end{bmatrix}.
    \label{eq:mdf-local-row}
\end{equation}
If the local velocity correction on this stencil is
\(\delta\vec u_a\), the marker correction is
\begin{equation}
    \mat J_a\delta\vec u_a = s_a .
    \label{eq:mdf-local-constraint}
\end{equation}
This equation is underdetermined because a single marker velocity constrains
several Eulerian stencil values. The term pseudoinverse means that, instead of
choosing one arbitrary inverse, we choose the inverse that gives the smallest
local correction on the stencil. For a single row vector, this
minimum-norm correction is obtained from the Moore--Penrose pseudoinverse,
\begin{equation}
    \delta\vec u_a
    = \mat J_a^+ s_a,
    \qquad
    \mat J_a^+
    = \mat J_a^\transp
    \left(\mat J_a\mat J_a^\transp\right)^{-1}
    = \frac{\mat J_a^\transp}{\sum_{i\in I_a}W_{ai}^2}.
    \label{eq:mdf-pseudoinverse}
\end{equation}
Equivalently, the target correction distributed over the stencil is
\begin{equation}
    \delta u_i
    = \frac{W_{ai}}{\sum_{j\in I_a}W_{aj}^2}\,s_a,
    \qquad i\in I_a .
    \label{eq:mdf-local-minimum-norm-update}
\end{equation}

In the immersed-boundary update this velocity correction is not applied
directly. It is generated by spreading a marker force-density increment through
the fluid mass update. If the local density and cell volume are approximated by
\(\rho_a\) and \(\Delta V_a\), then the force increment \(\delta\bm\Lambda_a\)
produces
\begin{equation}
    \delta u_i
    = -\Delta t\,\frac{q_a}{\rho_a\Delta V_a}
    W_{ai}\,\delta\Lambda_a .
    \label{eq:mdf-local-force-to-velocity}
\end{equation}
Matching \eqref{eq:mdf-local-force-to-velocity} with the pseudoinverse
correction \eqref{eq:mdf-local-minimum-norm-update} gives
\begin{equation}
    \delta\bm\Lambda_a
    =
    -\frac{\rho_a\Delta V_a}
    {\Delta t\,q_a\sum_{i\in I_a}W_{ai}^2}
    \vec s_a .
    \label{eq:mdf-local-force-increment}
\end{equation}
This is the local pseudoinverse used by MDF. It is also the inverse of the
diagonal self-mobility approximation
\begin{equation}
    D_a
    =
    \Delta t\,\frac{q_a}{\rho_a\Delta V_a}
    \sum_{i\in I_a}W_{ai}^2,
    \qquad
    \delta\bm\Lambda_a = -D_a^{-1}\vec s_a .
    \label{eq:mdf-diagonal-mobility}
\end{equation}
In the implementation, \(q_a\) is stored as the marker weight, \(\Delta V_a\)
is the local cell measure associated with the marker level, and
\(\sum_i W_{ai}^2\) is stored as the scalar \texttt{sumw2}. A relaxation
parameter \(\omega\) is included, so the computed increment is
\begin{equation}
    \delta\bm\Lambda_a^{(k)}
    =
    -\omega\,
    \frac{\rho_a\Delta V_a}
    {\Delta t\,q_a\sum_{i\in I_a}W_{ai}^2}
    \vec s_a^{(k)} .
    \label{eq:mdf-relaxed-force-increment}
\end{equation}
The accumulated marker force density after \(N_{\rm MDF}\) corrections is then
\begin{equation}
    \bm\Lambda_{\rm MDF}
    =
    \sum_{k=0}^{N_{\rm MDF}-1}\delta\bm\Lambda^{(k)} .
    \label{eq:mdf-accumulated-force}
\end{equation}

Written as an iteration for the slip, MDF has the form
\begin{equation}
    \vec s^{(k+1)}
    =
    \left(\mat I-\omega\mat A\mat D^{-1}\right)\vec s^{(k)},
    \label{eq:mdf-richardson-slip}
\end{equation}
where \(\mat D\) is the block-diagonal matrix with entries
\eqref{eq:mdf-diagonal-mobility}. If the markers are well separated relative
to the kernel support, \(\mat D\) is a good approximation to \(\mat A\) and a
single correction removes most of the slip. When nearby markers interact
strongly through overlapping kernels, the off-diagonal entries of \(\mat A\)
become important, and repeated MDF corrections reduce the remaining residual
without assembling or solving the full marker system.

\subsection{Validation}

\subsubsection{Two-Dimensional Validation}

\paragraph{Flow over a Stationairy Cylinder}

\subsubsection{Three-Dimensional Validation}

\paragraph{Flow over a fixed isolated sphere}

To validate the DF-IBM method in three dimensions, we consider the case of
uniform flow past a fixed sphere, where we may interpret the immersed boundary
force required to enforce no-slip as the hydrodynamic force on the sphere. We
may then compare the numerical drag coefficient to the Schiller-Naumann
correlation, which is typically reported \cite[84]{michaelides_multiphase_2022}
\begin{equation}
    C_D(\Reynolds_p) = \frac{24}{\Reynolds_p} \left(1 + 0.15 \Reynolds_p^{0.687}\right),
\end{equation}
and for which we may compare our scheme for the values \(\Reynolds_p \in \{10,50,100,200,250\}\).

For this case, we use a cubic fluid domain \(L_x = L_y = L_z = 40D\) with a
sphere of diameter \(D\) fixed at \(x_c = 5D\) from the inlet and centered in
\(y\) and \(z\) domain. The imposed far-field flow is
\begin{gather*}
    \vec u_{\infty} = U_\infty \vec e_x \qquad \text{at inlet,} \\
    \pdv*{\vec u}{x} = 0,\quad  p = 0, \qquad \text{at outlet.}
\end{gather*}
Then the particle Reynolds number is
\begin{equation*}
    \Reynolds_p = \frac{\rho_f U_\infty D}{\mu_f}.
\end{equation*}

Next we discuss the Lagrangian markers on the sphere; unlike the case for a
disc, the problem of choosing a 'uniform' distribution of makers remains an
open problem\autocite{saff_distributing_1997}; even the choice of what
constitutes 'uniform' is unclear. In this experiment, we distribute the points
according to a "fibbonaci spiral" algorithm:

% TODO: Describe the sphere algorithm used in RigidBodySphere.cpp

\begin{figure}[H]
    \centering
    \def\svgwidth{0.5\textwidth}
    \import{figure/vector}{dfibm-schiller-naumann-validation-domain.pdf_tex}
    \caption{Fluid domain and flow configuration for the flow over the isolated sphere. The particle is placed \(5\) particle diameters from the inlet where uniform inflow velocity \(\vec u_\infty\) is applied.}
    \label{fig:mdfibm-schiller-naumann-validation-domain-setup}
\end{figure}

\begin{figure}[H]
    \centering
    \import{figure/vector}{static-sphere-Re-sweep.tex}
    \caption{With \(\Delta t = 0.001\) and \(d_p / \Delta x = 25.6\) at the finest level.}
    \label{fig:mdfibm-schiller-naumann-validation-Re-sweep}
\end{figure}

% \printbibliography
